\subsection{Algebraic vs Topological smoothing}

\DMGtodo[inline]{algebraic and topological version of theorem axioms}

In this section we relate to notions of smoothing.

\textbf{Algebraic intersection conditions}

We say $f \colon \pi_1(S) \to \mathbb{R}$ is a \emph{class invariant function}, if $f(a)=f(b)$ whenever $a$ is conjugate to $b$. 
Note that a class invariant function is equivalent to a curve functional $f \colon \Curves(S) \to \mathbb{R}$.
A class invariant function $f$ satisfying the following conditions
for every $a,b \in \pi_1(S)$ hyperbolic is said to satisfy the \emph{algebraic intersection axioms}:

\begin{enumerate}
\item inverse invariant: $f(a)=f(A)$, for $A$ the inverse of $a$. 
\item stable: $f(a^n) = nf(a)$.
\item disconnected smoothing: For any pair of $R$-crossing elements $(a,b)$, 
 \begin{align}
    \label{eq:conn_oriented} f(a) + f(b) &\geq f(ab)
  \end{align}
\item connected smoothing: For any pair of $R$-parallel elements $(a,b)$, 
 \begin{align}
    \label{eq:conn_oriented} \mbox{oriented} \quad f(ab) &\geq  f(a) + f(b)\\
    \label{eq:conn_unoriented} \mbox{unoriented} \quad f(ab) &\geq f(aB).
  \end{align}
\end{enumerate}


\textbf{Topological intersection axioms}

Let $f \colon \Curves(S) \to \mathbb{R}$ be a curve functional satisfying

\begin{enumerate}
\item symmetric: $f([a]) = f([a^{-1}])$ for $a \in \pi_1(S)$.
\item stable: $f(C^n) = f(nC)$.
\item homogeneous: $f(nC)=nf(C)$
\item smoothing property:
     \[
      f\left(\mfig{curves-0}\right) \ge \max \left\{ f\left(\mfig{curves-1}\right), f\left(\mfig{curves-2}\right)\right\},
   \]
\end{enumerate}
We say that $f$ satisfies the \emph{topological intersection axioms}.

In this section we prove the following theorem.

\begin{theorem}
Let $f \colon \Curves(S) \to \mathbb{R}$ be a curve functional. $f$ satisfies the algebraic intersection axioms if and only if it satisfies the topological intersection axioms.
\end{theorem}

 Our main theorem is phrased in terms of topological smoothing, even though our proof uses the algebraic smoothing axioms. Both versions of the axioms are of independent interest, since the topological version is easier to verify for curve functionals arising from geometric functionals on the surface (such as length functions or intersection numbers) whereas the algebraic axioms are easier to check for curve functionals of more algebraic nature (such as periods of cross-ratios, stable lengths of word metrics, etc.).

\subsubsection{Triality}

To give a more workable and algebraic criterion for essential
crossings, we
adopt a group-theoretic point of view. For $a,b \in \pi_1(S,*)$, write
$[a]$ for the free homotopy class of curves represented by $a$;
algebraically this is the union of the conjugacy classes of~$a$ and
of~$A$. Write $[a][b]$ for the multi-curve $[a] \cup [b]$. Say that a
factorization $[ab]$ \emph{reduces} if some (equiv.\ every) taut
representative of $[ab]$ has a crossing so that one smoothing gives
$[a][b]$ and the other smoothing gives $[aB]$. Similarly say that
$[a][b]$ \emph{reduces} if some/every taut representative has a crossing
whose smoothings give $[ab]$ and $[aB]$. (This is an abuse of
notation, since the reduction is not invariant under individual
conjugation of $a$ and $b$.)

For $a,b \in \pi_1(S,*)$, let $a * b$ be a 4-valent graph mapped into
$S$ obtained by taking the union of the paths $a,b$ joined at the
basepoint~$*$, and let $[a * b]$
be its free homotopy class.
Free homotopy means in particular invariance under simultaneous
conjugation: $[a * b] = [caC * cbC]$. A \emph{minimal} representative of $[a * b]$
is one with a minimum number of transverse crossings.

Finally, if we fix a hyperbolic structure on~$S$, a non-trivial
element $a \in \pi_1(S,*)$ determines a geodesic $\ora{a}$ in
$\wt{S}$, the unique geodesic that stays a bounded distance from
$a^k \cdot \wt{*}$, with endpoints $a^+$ and $a^-$. Depending on the
cyclic order of the endpoints up to orientation reversal, we say that
$\ora{a}$ and $\ora{b}$ are
\begin{itemize}
\item \emph{parallel} if they appear in order $a^-,a^+,b^+,b^-$;
\item \emph{anti-parallel} if they appear in order $a^-,a^+,b^-,b^+$;
  and
\item \emph{cross} if they appear in order $a^-,b^-,a^+,b^+$.
\end{itemize}

\DMGtodo[inline]{definition of common powers}

\begin{definition}
  Non-trivial group elements $a,b\in \pi_1$ are \emph{common powers}
  if there are non-zero $k,\ell$ so that $a^k = b^\ell$. Since $\pi_1$
  has no torsion, we may assume that $k,\ell$ are relatively prime and
  then we see this is equivalent to the existence of a $c$ so that
  $a = c^\ell$ and $b = c^k$.
\end{definition}


\begin{proposition}\label{prop:crossings-triality}
  Let $a,b \in \pi_1(S,*)$ be non-trivial elements that are not
  common powers. Then in the following three trialities the same
  possibility holds in each case.
  \begin{enumerate}
  \item (Reduction) One of the following factorizations reduces
    to the other two:\\
    \begin{enumerate*}[itemjoin=\qquad]
    \item $[ab]$;
    \item $[aB]$; or
    \item $[a][b]$.
    \end{enumerate*}
  \item (Geometric) For every taut embedding of $a * b$, around the
    4-valent vertex we see in cyclic
    order, up to orientation reversal:\\
    \begin{enumerate*}[itemjoin=\qquad]
    \item $\mfig{curves-30}$;
    \item $\mfig{curves-31}$; or
    \item $\mfig{curves-32}$.
    \end{enumerate*}
  \item (Algebraic) The geodesics $\ora{a}$ and $\ora{b}$ are:\\
    \begin{enumerate*}[itemjoin=\qquad]
    \item parallel;
    \item anti-parallel; or
    \item cross.
    \end{enumerate*}
  \end{enumerate}
\end{proposition}

\begin{proof}
  Since $a$ and $b$ are not common powers, the endpoints of
  $\ora{a}$ and $\ora{b}$ are all distinct from each other (XXX: cite), so
  the algebraic triality is mutually exclusive and exhaustive. For the
  geometric triality, for any given embedding of $a * b$, one of the
  three given possibilities holds, but it is not \emph{a priori} clear that
  the same possibility holds for all taut embeddings. For the
  reduction triality, the possibilities are mutually exclusive (we
  can't have circular reductions, since at each reduction length with
  respect any Riemannian metric strictly decreases) but now it is not
  clear that any of the three possibilities holds.
  
  We start with the geometric triality. Consider any minimal representative
  $\rho$ of $[a * b]$ in one of the cases (a), (b), or~(c). Consider
  the (concrete) multi-curve~$\gamma$ given by replacing the
  $4$-valent vertex with a crossing. If we pick orientations on $\gamma$
  correctly, this is an embedding of $[ab]$, $[aB]$, or $[a][b]$,
  respectively, with a distinguished crossing~$x$.

  We claim that $\gamma$ is minimal. Otherwise, there is some reducing
  isotopy using RIII moves and some RI/RII reduction
  \cite{HS94:ShorteningCurves}. Any RIII move on~$\gamma$
  can be copied with an isotopy of~$\rho$, as can an RI or RII reduction of
  $\gamma$ that does not involve~$x$. An RI
  reduction at~$x$ is only possible if $a$ or $b$ is the identity, and
  an RII reduction involving~$x$ also reduces $\rho$,
  reducing the number of intersections of $\rho$ by one. In any case
  this contradicts minimality of~$\rho$ or
  non-triviality of $a,b$.

  We have thus shown that one of $[ab]$, $[aB]$, or $[a][b]$ has a
  minimal representative $\gamma$ with a crossing~$x$ that can be
  smoothed to give the other two multi-curves. Since we cannot have a
  circular chain of essential smoothings (for instance because we get
  strict inequalities of hyperbolic length), it follows that the cases
  are exclusive and thus that \emph{every} minimal representative of
  $[a * b]$ has the same cyclic order around the 4-valent vertex.
  Combined this
  proves equivalence of the first two trialities.

  For the remaining equivalence, suppose first we
  are in case~(c) of the geometric triality. For the fixed hyperbolic
  structure, take $\gamma$ to be the geodesic representative of the free
  homotopy class $[a][b]$, intersecting at the distinguished point
  $x \in S$, and let $\ora{a}$ and $\ora{b}$ be the respective
  lifts passing through some lift $\wt{x}$ of~$x$. Then $\ora{a}$ and
  $\ora{b}$ intersect only once and thus cross at infinity, so we are
  also in case~(c) of the algebraic triality.

  If we are in case~(a) of the geometric triality, again take the
  geodesic representative $\gamma$ of $ab$, with distinguished
  crossing~$x$. For concreteness, if necessary reverse the
  orientation of~$S$ so the branches appear in the order shown below.
  If we
  perform the disconnected (oriented) smoothing at $x$, we get two
  components $\gamma_a \in [a]$ and $\gamma_b \in [b]$; these concrete
  curves are typically not minimal. However, $\gamma_a$ and $\gamma_b$ are
  regular isotopic to a taut curve \cite[Lemma 5.5]{Thurston14:PBI},
  i.e., there is no RI move required in making them taut. This is
  equivalent to saying that their lifts $\wt{\gamma_a}$ and $\wt{\gamma_b}$ to
  $\wt{S}$ passing through a neighborhood of $\wt{x}$ are embedded in
  the disk (XXX: cite). Now consider
  the endpoints $a^+, a^-$ of $\wt{\gamma_a}$. As an oriented curve,
  $\wt{\gamma_a}$ is geodesic outside of neighborhoods of lifts of $x$, and
  inside those neighborhoods turns to the left. It then follows that
  $(ab)^+, a^+, a^-, (ba)^-$ occur in that cyclic order on
  $\partial \wt{S}$, and similarly $(ab)^-, b^-,b^+,(ba)^+$ occur
  in that cyclic order:
  \[
    \resizebox{40mm}{!}{\svgfig{lifts}}
  \]
  It therefore
  follows that $\ora{a}$ and
  $\ora{b}$ are parallel, so we are in case~(a) of the algebraic
  triality.

  Case~(b) follows in the same way, interchanging $b \leftrightarrow B$.
\end{proof}